\lecture{8}
\title{Classification and Logistic Regression}

\begin{document}

\subsection{Classification}

Over the past few lectures on linear regression,
we've assumed that the output space $\mathcal{Y}$ is $\mathbb{R}$.
In classification problems,
we instead take $\mathcal{Y}$ to be discrete;
for now, let's further assume $\mathcal{Y}$ is binary,
meaning $\mathcal{Y} = \{+1, -1\}$.

This is the setup for logistic regression,
covered in \href{/18.650/lectures/25/}{Lecture 25} of 18.650.
These notes will assume knowledge of that lecture.

\subsection{Bayes Classifier}

Here's the most common choice of loss function
for (binary) classification problems.

\textbf{Definition.}
The \textbf{zero-one loss} is
$\mathcal{L}(\hat{y}, y) := \mathds{1}[\hat{y} \neq y]$.
The loss equals zero if the classification was right,
and one otherwise.

In this case, it's clear that the best classifier
is the \textbf{Bayes classifier}.

\begin{theorem}{Bayes Classifier}
    The best (lowest-risk) classifier under zero-one loss is:
    \[ h(x) := \begin{cases}
    +1 & \text{ if } \mathbb{P}(Y = +1 \mid X = x) > \frac{1}{2}. \\
    -1 & \text{ otherwise}.
    \end{cases} \]
\end{theorem}

Of course, the challenge is that
we don't know what the distribution $[Y \mid X = x]$
exactly looks like.

\subsection{Nearest Neighbor Classifier}

Here's a simple design of a binary classifier.

\textbf{Definition.}
Given an input space $\mathcal{X}$, impose some distance metric
$d \colon \mathcal{X} \times \mathcal{X} \to \mathbb{R}^{\geq 0}$.
Then predict the label $h(x)$ like so:
\begin{quote}
    Identify the set $N_k(x)$
    of the indices of the $k$ training points closest to $x$.
    Then return $h(x) = \mathrm{sgn}\left(\sum_{i \in N_k(x)} y^{(i)}\right)$.
\end{quote}
This is called the \textbf{nearest neighbor classifier}
(or \textbf{NN}, or \textbf{KNN}).
It's simple, but slow:
making a single prediction requires
scanning the entire training set.

\subsection{ERM: Linear Classifier}

Let's try an ERM approach instead;
here's the set of hypotheses we'll restrict ourselves to.

\textbf{Definition.}
A \textbf{linear classifier} is a map
$h \colon \mathcal{X} \to \{+1, -1\}$ via
$h(x) := \mathrm{sgn}\left(w^{\top}x\right)$,
for some parameter $w$.

\begin{remark}
    We'll often pad every vector in $\mathcal{X}$ with an extra one, as usual;
    recall \href{/18.650/lectures/24/#affine-linear-regression}{Lecture 24}
    of 18.650.
\end{remark}

Then ERM tells us we should choose $w$
to minimize the empirical risk $\widehat{\mathrm{risk}}(h_w)$, given below:
\[ \widehat{\mathrm{risk}}(h_w)
:= \frac{1}{n} \sum_{i = 1}^n \mathcal{L}\left(w^{\top}x^{(i)} \cdot y^{(i)}\right)
~~~ \text{ where } ~~~
\mathcal{L}(z) := \mathds{1}[z < 0]. \]
Note that $\mathcal{L}(z)$ is just
a one-variable form of zero-one loss.
Unfortunately, the empirical risk described above
is hard to optimize (via gradient descent, say)
because $\mathcal{L}(z)$ is not convex or differentiable.
Here's a better alternative:

\textbf{Definition.}
A convex variant of zero-one loss is
\textbf{logistic loss}, given by
$\mathcal{L}_{\mathrm{logistic}}(z) := \log_2(1 + e^{-z})$.

\begin{center}
    \frame{\includegraphics[width=0.3\textwidth]{lecture-08/log-loss}}
\end{center}

\subsection{Addendum: Convexity}

We provide a brief reminder of convexity here;
it's important to know for optimization purposes.

\textbf{Definition.}
A twice-differentiable function
$f \colon \mathbb{R}^n \to \mathbb{R}$ is \textbf{convex}
if any of the following (equivalent) statements is true.
\begin{enumerate}
    \item We have $\nabla^2 f(x) \succeq 0$ for all $x$;
    that is, $\nabla^2 f(x)$ is positive semidefinite.
    \begin{itemize}
        \item Note that $\nabla^2 f(x)$ here is the \textbf{Hessian},
        which is the $n \times n$ matrix given by
        $\left[\nabla^2 f(x)\right]_{i, j} :=
        \frac{\partial^2 f}{\partial x_i \partial x_j}(x)$.
    \end{itemize}
    \item We have $f\left(\lambda \cdot u + (1 - \lambda) \cdot v\right)
    \leq \lambda \cdot f(u) + (1 - \lambda) \cdot f(v)$
    for all $u, v \in \mathbb{R}^n$ and $\lambda \in [0, 1]$.
    \begin{itemize}
        \item That is, any line segment connecting two points
        on the graph of $f$ lies above the graph.
    \end{itemize}
    \item We have $f(u) \geq f(v) + \nabla f(v)^{\top}(u - v)$
    for all $u, v \in \mathbb{R}^n$.
    \begin{itemize}
        \item That is, the tangent plane to $f$ at $v$
        lives underneath the graph of $f$ everywhere.
    \end{itemize}
\end{enumerate}

Here's why we care about convexity so much:

\begin{theorem}{Local = Global}
    For any convex objective over a convex domain,
    any local minimum is a global minimum.
\end{theorem}

So techniques that seek out local minima (e.g., gradient descent)
are really seeking out the global minimum.
You can check that logistic loss is convex.

\subsection{MLE: Logistic Regression}

As in \href{/18.650/lectures/25/}{Lecture 25},
we'll switch our hypothesis to be a map
$p_w \colon \mathcal{X} \to [0, 1]$,
where the goal is to have
$p_w(x)$ be a solid approximation of $\mathbb{P}(Y = +1 \mid X = x)$.
Then our classification model is
$h_w(x) := \mathrm{sgn}\left(p_w(x) - \frac{1}{2}\right)$, naturally.

Specifically, we'll assume that $p_w$ takes the form
$p_w(x) = \sigma\left(w^{\top}x\right)$,
where $\sigma(z) := \frac{e^z}{1 + e^z}$
is the sigmoid function.

\begin{theorem}{Logistic Loss}
    The log (base 2) likelihood of $p_w(x)$ is:
    \[ \ell(w) = -\sum_{i = 1}^n
        \log_2\left(
            1 + \exp\left(-y^{(i)}w^{\top}x^{(i)}\right)
        \right)
    = - \sum_{i = 1}^n \mathcal{L}_{\mathrm{logistic}}\left(
        w^{\top} x^{(i)} \cdot y^{(i)}
    \right). \]
    The $w^{\top}x^{(i)} \cdot y^{(i)}$ term is called
    the \textbf{margin}, whose sign determines correctness
    and whose magnitude determines confidence.
\end{theorem}

The proof is just plain computation, so we omit it.
Note that this differs slightly from the result in
\href{/18.650/lectures/25/#mle-for-logistic-regression}{Lecture 25}
of 18.650 because of the convention that
$y^{(i)} \in \{+1, -1\}$ instead of
$y^{(i)} \in \{0, 1\}$.

The interesting part is the coincidence that
logistic loss shows up once again.
The following are therefore identical:
\begin{itemize}
    \item Maximizing the likelihood, given a model
    $p_w(x) = \sigma(w^{\top}x)$ designed to approximate
    $\mathbb{P}(Y = +1 \mid X = x)$.
    \item Minimizing the logistic loss, given a hypothesis
    $h(x) := \mathrm{sgn}(w^{\top}x)$ designed to
    predict the classification of $x$.
\end{itemize}

\subsection{Addendum: Log-Odds and Linear-Logistic Equivalence}

\textbf{Definition.}
The \textbf{log-odds} is defined as the value
$f(x) := \log \left[
    \frac{\mathbb{P}(Y = +1 \mid X = x)}{\mathbb{P}(Y = -1 \mid X = x)}
\right]$.

\begin{theorem}{Linear-Logistic Equivalence}
    Given a classification model,
    consider the following two procedures.
    \begin{itemize}
        \item Logistic regression,
        seeking a parameter $w_1$ for which
        $p_{w_1}(x) := \sigma(w_1^{\top}x)$
        equals the probability
        $\mathbb{P}(Y = +1 \mid X = x)$.
        \item Linear regression,
        seeking a parameter $w_2$ for which
        $h_{w_2}(x) := w_2^{\top}x$
        equals the log-odds $f(x)$.
    \end{itemize}
    Then the two procedures seek the same parameter: $w_1 = w_2$.
\end{theorem}

\begin{proof}
    Obvious; the log-odds function $f(x)$
    is just the inverse of $\sigma(\dots)$, in that
    $\log\left[
        \frac{\sigma(w_1^{\top}x)}{1 - \sigma(w_1^{\top}x)}
    \right] = w_1^{\top}x$.
    ~ $\blacksquare$
\end{proof}

So all of the discussion about features and kernels
concerning linear regression
applies equally well to logistic regression.

\end{document}